\documentclass[../DoAn.tex]{subfiles}
\begin{document}
\section{Kết luận và hướng phát triển}
\subsection{Kết luận}
Trong những năm gần đây, nhiều hệ thống và phần mềm quản lý tài sản công nghệ thông tin đã được phát triển và đưa vào sử dụng, từ các giải pháp thương mại cho đến các nền tảng mã nguồn mở như Ứng dụng quản lý tài sản. So với các sản phẩm thương mại, hệ thống quản lý tài sản được xây dựng trong đồ án có ưu điểm là chi phí triển khai thấp, khả năng tùy biến cao và phù hợp với nhu cầu của các tổ chức có quy mô vừa và nhỏ. So với các nghiên cứu và ứng dụng tương tự dựa trên Ứng dụng quản lý tài sản, đồ án không chỉ dừng lại ở việc cài đặt và sử dụng hệ thống có sẵn, mà còn tập trung vào việc phân tích, tổ chức lại kiến trúc, thiết kế mô hình dữ liệu và triển khai hệ thống trong một môi trường thực tế cụ thể.

Trong suốt quá trình thực hiện đồ án tốt nghiệp, sinh viên đã từng bước hoàn thành các nội dung chính từ phân tích yêu cầu, thiết kế hệ thống, xây dựng ứng dụng cho đến kiểm thử và triển khai thử nghiệm. Cụ thể, sinh viên đã xây dựng được một hệ thống quản lý tài sản hoạt động ổn định, đáp ứng các nghiệp vụ cơ bản như quản lý thông tin tài sản, quản lý người dùng, cấp phát và thu hồi tài sản, cũng như theo dõi lịch sử sử dụng tài sản. Bên cạnh đó, đồ án còn hoàn thiện bộ tài liệu kỹ thuật tương đối đầy đủ, phục vụ cho việc bảo trì và mở rộng hệ thống trong tương lai.

Tuy nhiên, do giới hạn về thời gian và phạm vi của một đồ án tốt nghiệp, hệ thống vẫn còn một số hạn chế nhất định. Một số chức năng nâng cao của Ứng dụng quản lý tài sản chưa được khai thác triệt để, khả năng tùy biến giao diện và tối ưu trải nghiệm người dùng còn ở mức cơ bản. Việc đánh giá hiệu năng và khả năng chịu tải của hệ thống mới chỉ dừng lại ở mức thử nghiệm với quy mô nhỏ, chưa phản ánh đầy đủ các kịch bản sử dụng trong môi trường có số lượng người dùng lớn.

Những đóng góp nổi bật của đồ án có thể được tổng hợp ở việc đề xuất và hiện thực hóa giải pháp tổ chức, tùy biến một nền tảng mã nguồn mở để xây dựng hệ thống quản lý tài sản phù hợp với nhu cầu thực tế. Thông qua đồ án, sinh viên đã tích lũy được nhiều kinh nghiệm quan trọng trong việc phân tích bài toán thực tế, thiết kế kiến trúc hệ thống, làm việc với cơ sở dữ liệu, cũng như triển khai và vận hành một ứng dụng web trong môi trường máy chủ. Đây là những bài học có giá trị, tạo nền tảng cho việc nghiên cứu và phát triển các hệ thống phần mềm phức tạp hơn trong tương lai.
\subsection{Hướng phát triển}

Trong thời gian tới, để hoàn thiện hơn hệ thống quản lý tài sản đã xây dựng, cần tiếp tục thực hiện một số công việc nhằm nâng cao chất lượng và tính ứng dụng của sản phẩm. Trước hết, các chức năng hiện có cần được rà soát và tối ưu, đặc biệt là giao diện người dùng và quy trình thao tác, nhằm nâng cao trải nghiệm sử dụng và giảm thiểu các thao tác không cần thiết trong quá trình quản lý tài sản.

Bên cạnh đó, hệ thống có thể được mở rộng bằng cách bổ sung các chức năng nâng cao như phân quyền chi tiết theo vai trò, tích hợp cơ chế thông báo tự động khi tài sản sắp đến hạn bảo trì hoặc hết thời gian sử dụng, cũng như hỗ trợ xuất báo cáo thống kê dưới nhiều định dạng khác nhau. Việc khai thác sâu hơn các API của Ứng dụng quản lý tài sản cũng là một hướng đi tiềm năng, giúp hệ thống linh hoạt hơn trong việc tích hợp với các ứng dụng khác trong tổ chức.

Về mặt triển khai, trong tương lai hệ thống có thể được thử nghiệm với quy mô lớn hơn, trên hạ tầng có cấu hình cao hơn hoặc trên các nền tảng điện toán đám mây. Điều này cho phép đánh giá toàn diện hơn về hiệu năng, khả năng chịu tải và độ ổn định của hệ thống trong các kịch bản sử dụng thực tế. Ngoài ra, việc bổ sung các công cụ giám sát và ghi log cũng sẽ giúp nâng cao khả năng vận hành và bảo trì hệ thống.

Xa hơn, các hướng nghiên cứu mới có thể tập trung vào việc ứng dụng các kỹ thuật phân tích dữ liệu để khai thác thông tin từ lịch sử sử dụng tài sản, phục vụ cho việc ra quyết định quản lý. Đây là hướng phát triển tiềm năng, góp phần nâng cao giá trị thực tiễn của hệ thống quản lý tài sản trong bối cảnh chuyển đổi số hiện nay.

\end{document}